% paper/sections/13f_error_budget.tex
%
% V1-V11 横断の精度バジェット総合
% Hosts labels: sec:error_budget, sec:accuracy_summary
% (tab:accuracy_summary 09f, sec:time_accuracy_table 07 を authoritative とし
%  本節からは hosting せず外部参照のみで連結する)
% ═══════════════════════════════════════════════════════════════════════════════
\subsection{精度バジェット：V1--V11 誤差寄与の統合評価}
\label{sec:error_budget}
\label{sec:accuracy_summary}

\paragraph{構成．}
本節は V1--V11 横断の精度結果を一表に集約し，
\textbf{各誤差源} (空間離散，時間離散，界面表現，曲率推定，CLS 再初期化，射影後面速度，履歴圧力) と
\textbf{各設定パラメータ} ($\rho_r$, $\alpha$, $\Delta t$, $h$) の関係を整理する
（観測収束次数の Richardson 外挿運用は~\cite{Roy2003}）．
本節は，最終スキームについて，
どの誤差源がどの検証問題で支配的になるかを整理する精度バジェットである．
U10/U11 は第~\ref{sec:numerical_integrity}章の単体条件として扱い，
U12/V11 はアクティブ幾何毛管分解の統合前ゲートとして扱う．
保存形共通フラックスの有限振幅物理応答と周期壁上昇気泡応答は，
第~\ref{sec:benchmarks}章のベンチマークと事前適合性確認で扱う．

\paragraph{判定基準（§\ref{sec:verification} Type 分類との整合）．}
本節は §\ref{sec:verification}「判定区分と Type 分類」に基づき
Type-A/Type-B/Type-D の三類型で各 V テストの判定を集約する
（凡例は表~\ref{tab:V_verdict_types}）．
V1-a, V1-b, V2, V4, V5 は \textbf{Type-A 判定基準}：
連続方程式または単体要素の理想化次数ではなく，
実際に測定する縮約離散恒等関係に対する文献根拠付き判定基準で評価する．
V1-a は単相 TGV 空間の FVM PPE 結合系に対する $O(h^2)$ 傾向，
V1-b は非増分の標準 PPE 射影の $O(\Delta t)$ 傾向，
V2 は壁片側 Kovasznay 統合残差の $O(h^4)$ 傾向，
V4 は固定壁一様オフセット診断の $O(\|\bU\|)$ 有界残差，
V5 は CCD 寄生流れ絶対値 $u_\infty^{\mathrm{end}}$ の $\sigma/\mathrm{Re}_c$ スケール基準
を合格基準とする．
V10-a 質量保存軸 / V10-b 質量保存軸は \textbf{Type-B}：
Olsson--Kreiss 質量補正~\cite{OlssonKreiss2005} を 10 step ごとに適用する
構造的アルゴリズム修正による合格化．
V7（毛管圧力ジャンプ/射影の界面帯に律速され $2.0$ に届かない結合系収束率），V10-a 形状軸（固定格子 CLS 位相/しきい値化誤差 + slot/h$=6.4$ 解像不足），
V10-b 形状軸（格子幅まで薄くなる folded filament の Eulerian 固定格子限界）は \textbf{Type-D}：
連続方程式・離散構造に由来する理論的下限と整合するものとして
条件付き合格 ($\triangle$)．これらは格子増強だけで合格化するのではなく，
norm 切替，項別誤差分解，dt/h 独立掃引，時系列解析により
保証範囲を別途定める必要がある．
V9 は Type-A/B/D の分類ではなく，圧力ジャンプ・毛管射影結合系と同時に基準/局所
$\varepsilon$ 幅を切り替えても有限時間静止液滴を破綻させないことを測る
\textbf{演算子群条件付き診断}として別枠で扱う．
V11 も Type-A/B/D ではなく，アクティブ幾何毛管分解を標準物理経路へ入れる前の
\textbf{採用境界ゲート}として別枠で扱う．

\paragraph{統合検証成績一覧．}

\begin{table}[!ht]
\caption{V1--V11：統合精度まとめ}
\label{tab:v_accuracy_summary}
\PaperCaptionNote{収束率は log-log fitting (3 点以上)；単一値はその設定での代表値．本表は Chapter 12 の成果総括表~\ref{tab:verification_summary} と同じ「ID / テスト / 期待・判定基準 / 実測・観測値 / 判定」文法で読む．表~\ref{tab:U_summary} は成果表ではなく U 系列の検証設計マップである．V10-a/V10-b は質量保存軸 (Type-B) と形状軸 (Type-D) に分離するため，表~\ref{tab:V_summary} の V10-a/b から 1:2 に展開される（V11 ゲートを含め計 15 行）．補助観測指標および全観測指標の網羅的列挙は表~\ref{tab:V_summary} を参照．}
\centering
\PaperTableExtraDenseSetup
\begin{tabularx}{\linewidth}{@{}>{\raggedright\arraybackslash}p{0.10\linewidth}>{\raggedright\arraybackslash}p{0.18\linewidth}>{\raggedright\arraybackslash}p{0.23\linewidth}>{\raggedright\arraybackslash}X>{\raggedright\arraybackslash}p{0.15\linewidth}@{}}
\toprule
ID & テスト & \shortstack{期待・判定\\基準} & \shortstack{実測・観測\\値} & 判定 \\
\midrule
V1-a & 単相 TGV 空間 & FVM PPE 結合系の $O(h^2)$ 傾向 & slope($E_{\mathrm{rel}}, h$) $=\mathbf{2.00}$；$E_{\mathrm{rel}}=1.583\times 10^{-7}$ at $N=64$ & \checkmark Type-A \\
V1-b & 単相 TGV 時間 & 標準 PPE の $O(\Delta t)$ 傾向 & slope($E_{\mathrm{rel}}, \Delta t$) $=\mathbf{1.00}$ & \checkmark Type-A \\
V2 & Kovasznay 空間残差 & 壁片側残差 $O(h^4)$ & slope($\|R\|_\infty, h$) $=\mathbf{3.95}$ & \checkmark Type-A \\
V3 & 静止液滴 long-term & BF 平衡 + 標準構成静止判定 & $|\Delta p_{\mathrm{rel}}|=0.97\%$ at $N=128$；標準構成 $N=32,T=10$: $K_{\max}=1.72\!\times\!10^{-6}$, $|\Delta V|/V_0\le2.92\!\times\!10^{-15}$, $D=0$ & \checkmark \\
V4 & 一様オフセット残差 & $O(\|\bU\|)$ 有界残差，運用上限 $<3$ & $\|\Delta u\|_\infty^{\max}/\|\bU\|=2.63$ & \checkmark Type-A \\
V5 & 寄生流れ 200 step & CCD $u_\infty^{\mathrm{end}}$ の $\sigma/\mathrm{Re}_c$ スケール基準 & $1.93\!\times\!10^{-2}$ at $\rho_r=1,N=32$；FD 比 $1/5.44$ 補助 & \checkmark Type-A \\
V6 & 密度比掃引 & ブローアップなし；CLS 体積ドリフトは丸め誤差床；毛管面バランス & 全 8 ケース安定；$u_\infty^{\mathrm{end}}\le1.3\!\times\!10^{-9}$；$|\Delta V_\psi|/V_{\psi,0}\le1.2\!\times\!10^{-16}$ & \checkmark \\
V7 & IMEX-BDF2 二相時間 & 毛管圧力ジャンプ/射影の界面帯制約下の結合系実効収束率 & slope($\|e\|_\infty,\Delta t$) $=\mathbf{1.59}$ & $\triangle$ Type-D \\
V8 & 非一様 NS 静止 & 非一様 BF 静止判定基準 & $|\Delta p_{\mathrm{rel}}|=2.86\%$ at $\alpha=2,N=96$ & \checkmark \\
V9 & 基準/局所 $\varepsilon$ 切替 & 圧力ジャンプ・毛管射影結合系同時使用下の有界性診断 & B/C=$1.00$ at $N=32$；$u_\infty^{\max}\le9.82\!\times\!10^{-10}$；$|\Delta V|/V_0\le4.73\!\times\!10^{-16}$ & \checkmark 結合系診断 \\
V10-a-mass & Zalesak 質量保存軸 & Olsson--Kreiss 質量補正後 drift $<0.5\%$ & $4.70\!\times\!10^{-4}\%$ at $N=128$ & \checkmark Type-B \\
V10-a-shape & Zalesak 形状軸 & 固定格子位相/しきい値化誤差 + slot/h$=6.4$ 解像不足と整合 & centroid err $4.911\!\times\!10^{-3}$；shape $L^1=8.307\!\times\!10^{-3}$ & $\triangle$ Type-D \\
V10-b-mass & 単一渦 reversal 質量保存軸 & Olsson--Kreiss 質量補正後 drift $<0.5\%$ & $0.0428\%$ at $N=128,T=8$ & \checkmark Type-B \\
V10-b-shape & 単一渦 reversal 形状軸 & 格子幅まで薄くなる folded filament の Eulerian 固定格子下限と整合 & $L^1$ reversal error $=2.248\times10^{-2}$ & $\triangle$ Type-D \\
V11 & アクティブ幾何毛管分解ゲート & 全圧力像への過射影を棄却し，$R_p$ 未確定の分解を標準経路へ入れない & 厳密全圧力分解は毛管波駆動 $0$；成分体積 Hodge 診断は N32/N64 $2.1176/2.3055$；未受理分解は標準経路から分離 & \checkmark ゲート \\
\bottomrule
\end{tabularx}
\end{table}

\FloatBarrier

\paragraph{誤差源の整理．}
\begin{enumerate}\setlength{\itemsep}{0pt}
  \item \textbf{空間離散誤差}．CCD 内点 6 次設計に対し，V2 Kovasznay residual は
        壁片側二階微分境界閉包と大域コンパクト結合により実効 $3.95$．
        Type-A 判定基準 $O(h^4)$ を満たす．TGV では時間誤差床下に到達するため空間収束率は飽和．
  \item \textbf{時間離散誤差}．単相 NS の AB2 predictor は単体 $O(\Delta t^2)$ 設計だが，
        V1-b は標準射影結合を測るため，合格基準は $O(\Delta t)$ 傾向である．
        二相結合では毛管圧力ジャンプ / アフィン FCCD 射影の界面帯誤差が
        結合系実効収束率 $1.59$（V7; \autoref{sec:imex_bdf2_twophase_time}）を支配し，
        純粋 BDF2 2 次ではなく毛管結合系の実効次数として扱う．
  \item \textbf{界面表現誤差}．Heaviside--$\delta$ の基準/局所 $\varepsilon$ 切替は
        FCCD, HFE, 圧力ジャンプ PPE, 毛管成分の圧力勾配射影，射影後面速度閉包，アフィン圧力履歴面加速度と同時に使う条件で有限時間安定（V9）；
        ただしこの direct-$\psi$/HFE 確認では B/C が同一で，局所幅を独立した改善機構とは主張しない．
        CCD/CSF 縮約経路の結果は圧力ジャンプ/HFE 結合系の精度根拠には用いない．
        さらに射影後面速度閉包を用いる結合系では，
        $\psi$ 移流の速度表現そのものが誤差源になりうる．
        射影で得た正準面速度を節点へ再構成してから面フラックスを作り直す経路は，
        式~\eqref{eq:projection_native_psi_transport} の保存則と同値ではないため，
        static-interface の $\psi$ 摂動を曲率誤差へ増幅する．
        同じ理由で，affine-jump IPC の履歴圧力項も
        不連続な $p^n$ を節点スカラーとして微分するのではなく，
        式~\eqref{eq:affine_pressure_history_faces} の正準面加速度として保持する必要がある．
  \item \textbf{射影残差誤差と圧力代表誤差の分離}．
        DC 反復回数そのものは精度証明ではなく，
        式~\eqref{eq:ppe_dc_residual_contract} の高次残差が射影受理条件である．
        したがって圧力ジャンプ PPE の精度評価では，
        楕円求解残差，投影後発散，圧力代表の選択を別々の観測量として扱う．
        残差受理条件を満たした後に残る界面帯圧力振動は，
        直ちに運動量圧力仕事の破綻とは解釈しない．
        affine-jump 経路の物理仕事は面共鎖
        $A_f(G_fp-B_f(j))$ で定義されるため，
        拡散界面帯の未補正節点圧力は Hodge 代表または相内代表に写してから読む．
        壁付き面速度状態空間については U11 が
        $C_wP_w=0$，計量自己随伴性，制限 Green 恒等式，および
        壁/周期壁商空間の階数条件を丸め誤差まで確認したため，
        V6/V7/V9 の圧力ジャンプ/HFE 統合診断はこの単体条件の上で読む．
        ただし U11 は小型代数条件であり，周期壁の有限振幅物理応答を
        第13章に追加するものではない．
  \item \textbf{アクティブ幾何毛管分解の圧力反力分解ゲート}．V11 は，
        式~\eqref{eq:capillary_pressure_reaction_projection} の $R_p$ を定めない
        全圧力像への分解が，毛管波の面上駆動を代数的に消し得ることを
        統合前に検出する．成分体積 Hodge 診断は非静的波の検出器として有用だが，
        それ自体は最終 $R_p$ の代替証明ではない．したがって $R_p$ 未確定の分解は
        未受理とし，標準圧力ジャンプ経路の V6/V7/V9 と第~\ref{sec:benchmarks}章の
        毛管波結果から分離する．
  \item \textbf{曲率推定誤差}．V5 では CSF $\kappa_{lg}$ + CCD 勾配の
        $u_\infty^{\mathrm{end}}$ 絶対値が全表ケースで $2\times10^{-2}$ 以下に保たれ，
        参照として FD 値を全条件で下回る．
        粗格子・低密度比ほど CCD/FD の差は大きい．
  \item \textbf{CLS 形状追跡誤差（質量保存軸 vs 形状精度軸の分離）}．V10 の Zalesak は
        $\alpha=1$ の一様格子診断に固定する．Olsson--Kreiss 質量補正~\cite{OlssonKreiss2005}
        を 10 step ごとに適用した結果，$N=128$ で質量ドリフト $4.70\times 10^{-4}\%$（合格基準
        $0.5\%$ を 3 桁下回り \checkmark）に到達する一方，centroid err $4.911\times 10^{-3}$,
        shape $L^1$ $8.307\times 10^{-3}$, 重心収束率 $\approx 1.22$ は固定格子
        CLS 位相/しきい値化誤差と slot/h $=6.4$ の sharp-corner 解像不足に
        律速される．単一渦 (V10-b) も同質量修正下で
        質量ドリフト $0.0428\%$ \checkmark / $L^1$ reversal $2.248\times10^{-2}$ ($T=8$ の
        格子幅 folded filament に対する Eulerian 固定一様格子の理論的下限~\cite{EnrightFedkiwFerzigerMitchell2002}) と
        同じく\textbf{保存軸 \checkmark / 形状軸 $\triangle$} の二軸構造を取る．
  \item \textbf{判定軸の分離}．V7/V10 系は単一の設計次数だけで判定せず，
        毛管/射影結合，固定格子 CLS 形状誤差，および
        質量補正後の保存軸を分けて読む．
\end{enumerate}

\paragraph{合格判定要約（§\ref{sec:verification} Type 分類で集約）．}
\begin{itemize}\setlength{\itemsep}{0pt}
  \item \checkmark \textbf{合格（設計基準）}:
        V3 (BF 平衡，$|\Delta p_{\mathrm{rel}}|=0.97\%$ at $N=128$；標準構成静止判定で $K_{\max}\le1.72\!\times\!10^{-6}$, $|\Delta V|/V_0\le2.92\!\times\!10^{-15}$, $D=0$)，
        V6 (密度比ロバスト, $\rho_r=833$ までブローアップなし; $u_\infty^{\mathrm{end}}\le1.3\!\times\!10^{-9}$; $|\Delta V_\psi|/V_{\psi,0} \le 1.2\!\times\!10^{-16}$)，
        V8 (非一様静止, $\alpha=2,N=96$ で $|\Delta p_{\mathrm{rel}}|=2.86\%$).
  \item \checkmark \textbf{Type-A 判定基準合格}:
        V1-a (TGV 空間 FVM PPE $O(h^2)$ 傾向, 細格子側局所収束率 $2.00$；補助: $E_{\mathrm{rel}}=1.583\!\times\!10^{-7}$ at $N=64$)，
        V1-b (標準射影結合 $O(\Delta t)$ 傾向, 収束率 $1.00$)，
        V2 (Kovasznay 壁残差 $O(h^4)$ 傾向, 収束率 $3.95$)，
        V4 (固定壁一様オフセット $O(\|\bU\|)$ 有界残差, $\|\Delta u\|_\infty^{\max}/\|\bU\| = 2.63$)，
        V5 (CCD 寄生流れ絶対値 $1.93\!\times\!10^{-2}$ at $\rho_r=1, N=32$).
  \item \checkmark \textbf{Type-B 構造修正 合格}:
        V10-a 質量保存軸 (Olsson--Kreiss 質量補正; $N=128$ 質量ドリフト $4.70\!\times\!10^{-4}\%$)，
        V10-b 質量保存軸 (同; $0.0428\%$).
  \item $\triangle$ \textbf{Type-D 理論的下限との整合}:
        V7 (二相結合系の実効収束率 $1.59$ — 毛管圧力ジャンプ/射影
        界面帯に律速され，BDF2 単体 $2.0$ には届かない),
        V10-a 形状軸 ($\alpha=1$ Zalesak centroid err $4.911\times 10^{-3}$,
        shape $L^1$ $8.307\times 10^{-3}$ — 固定格子 CLS 位相/しきい値化誤差と
        slot/h $=6.4$ sharp-corner 解像不足律速),
        V10-b 形状軸 ($T=8$ 単一渦 reversal $L^1=2.248\times10^{-2}$
        — 格子幅 folded filament の Eulerian 固定格子原理的下限~\cite{EnrightFedkiwFerzigerMitchell2002}).
  \item $\times$ \textbf{要改善 / 推奨せず}: 該当なし.
  \item \checkmark \textbf{演算子群条件付き診断}: V9
        (基準/局所 $\varepsilon$ 切替は $N=32$ で B/C=$1.00$，$u_\infty^{\max}\le9.82\!\times\!10^{-10}$，$|\Delta V|/V_0\le4.73\!\times\!10^{-16}$；
        HFE 曲率 + 圧力ジャンプ PPE + 毛管成分の圧力勾配射影 + 射影後面速度閉包 + アフィン圧力履歴面加速度と同時使用しても有限時間静止液滴を破綻させないが，
        局所幅の改善効果はこの縮約確認では主張しない).
\end{itemize}

\paragraph{設計指針 (PR-5 Algorithm Fidelity)．}
\begin{itemize}\setlength{\itemsep}{0pt}
  \item 単相 NS：CCD 内点 6 次空間 + AB2 predictor 2 次時間設計は単体レベルで維持する．
        統合縮約実験では，Kovasznay 壁残差は $O(h^4)$，
        標準 PPE TGV 時間傾向は $O(\Delta t)$ として判定する．
  \item 二相 NS：圧力ジャンプ・毛管射影結合系
        （分相 PPE + DC + HFE + 毛管成分の圧力勾配射影 + 射影後面速度閉包 + アフィン圧力履歴面加速度）が
        密度比 $\le 833$ でロバスト（V6）．
        ただし時間方向の実効次数は界面処理に律速されるため，
        \textbf{$\Delta t$ を $h_{\min}$ より小さく取る安全ファクタ}（CFL $\le 0.4$）が推奨．
  \item 非一様格子：$\alpha>1$ の NS 静止問題と基準/局所 $\varepsilon$ 切替は V8/V9 で検証し，
        Zalesak slot と単一渦 reversal は $\alpha=1$ の一様格子 CLS 移流診断に固定する．
        局所幅を改善根拠として用いるには，移動界面または再初期化が局所厚みを
        実際に経由する別検証を要求する．
  \item 寄生流れ抑制：CCD の $u_\infty^{\mathrm{end}}$ 絶対値は V5 全表ケースで
        $O(10^{-2})$ 以下に保たれ，代表的には
        $\rho_r=1,N=32$ で $1.93\!\times\!10^{-2}$，$\rho_r=1,N=128$ で $6.28\!\times\!10^{-4}$．
        参照として FD 値を全条件で下回り（$\rho_r=1,N=32$ で $1/5.44$，
        $\rho_r=100,N=64$ で $1/2.5$），高解像度・高密度比では CCD/FD 差が縮小する．
\end{itemize}

\paragraph{検証対応の統合．}
本節は V1--V11 全結果の横断統合点である．各実験が測定した式・離散化・物理指標の対応を
以下に集約する：
\begin{itemize}\setlength{\itemsep}{0pt}
  \item V1 (\autoref{sec:single_phase_ns}) — Navier--Stokes \eqref{eq:NS_onefluid}
        と AB2 + 標準 PPE 縮約経路に対するエネルギー減衰と非圧縮残差．
  \item V2 (\autoref{sec:single_phase_ns}, \autoref{sec:kovasznay}) —
        CCD 内点 6 次演算子と壁片側境界閉包を含む Kovasznay 統合残差．
  \item V3, V5 (\autoref{sec:bf_static_droplet}, \autoref{sec:spurious_multistep})
        — BF 力学的バランス \eqref{eq:bf_fccd_unification} + CSF
        \eqref{eq:csf_final} に対する Laplace 圧維持と寄生流れ，および
        圧力ジャンプ/Riesz 標準構成の静止判定による幾何・体積保持．
  \item V4 (\autoref{sec:galilean_offset}) —
        固定壁系の一様オフセット残差に対する $O(\|\bU\|)$ 有界性．
  \item V6, V7 (\autoref{sec:density_ratio_sweep},
        \autoref{sec:imex_bdf2_twophase_time}) — 圧力ジャンプ・毛管射影結合系
        （分相 PPE \eqref{eq:split_ppe} + DC + HFE + 毛管成分の圧力勾配射影 +
        射影後面速度閉包 + アフィン圧力履歴面加速度）に対する密度比安定性と二相時間精度．
  \item V8 (\autoref{sec:nonuniform_ns_static}) — 非一様 NS に対する静止液滴
        $\Delta p_{\mathrm{rel}}$ 維持．
  \item V9 (\autoref{sec:local_eps_validation}) — 圧力ジャンプ・毛管射影結合系上の基準/局所 $\varepsilon$ 切替に対する
        寄生流れ，体積分率ドリフト，および B/C 同一性．
  \item V10-a (\autoref{sec:zalesak_cls_advection}) — Zalesak CLS 移流に対する
        質量ドリフトと centroid 追跡精度（質量保存軸 / 形状軸の二軸構造）．
  \item V10-b (\autoref{sec:single_vortex_reversal}) — 単一渦 deformation
        reversal に対する $L^1$ 誤差と質量ドリフト（同二軸構造）．
  \item V11 (\autoref{sec:ao_fast_capillary_gate}) — アクティブ幾何毛管分解の
        全圧力像への過射影反例，成分体積 Hodge 診断，標準経路への採用境界．
\end{itemize}
本節 \autoref{tab:v_accuracy_summary} はこの全対応の精度結果を一表に集約する
唯一の総合表であり，他章 (\autoref{sec:conclusion}) からの参照は
本表をエンドポイントとする．

\FloatBarrier
